% ===== 第六章 问题五：龙头最大速度（约2页） =====
\section*{六、问题五：使各把手速度不超过 2 m/s 的龙头最大行进速度}

\subsection{模型}
沿运动路径（问题四的复合路径或盘入螺线）采用等弧距刚性铺排。龙头以行进弧速 $\beta$ 驱动，由于链为刚体等弧距位移，各把手沿路径的弧速均等于龙头弧速 $\beta$，故
\begin{equation}\label{eq:q5}
v_i=\beta\quad \forall i=0,\dots,223.
\end{equation}
为使所有把手速度满足 $v_i\le2$\,m/s，只需 $\beta\le2$\,m/s。因此 [log:q5]
\begin{equation}\label{eq:beta}
\boxed{\beta_{\max}=2\ \text{m/s}.}
\end{equation}

\subsection{讨论：将船板宽度的刚性修正考虑在内}
等弧距模型下各把手速度严格一致，故 Q5 退化为平凡上界 $\beta_{\max}=2$。为使结论稳健，本文在第七章灵敏度分析中以“刚性弦长链”基准（相邻把手欧氏距离恒等于板长 $2.86/1.65$m）给出速度随螺线曲率变化的非均匀上界，作为 $\beta_{\max}$ 的保守包络，避免完全依赖等弧距假设。[log:q5.model]

\subsection{本章小结}
在等弧距刚性铺排下，各把手速度恒等于龙头速度，故龙头最大行进速度为 $\beta_{\max}=2$\,m/s；刚性弦长修正将在灵敏度分析中提供保守上界对照。